%%%%%%%%%%%%
%
% $Autor: Wings $
% $Datum: 2019-12-09 11:47:41Z $
% $Pfad: komponenten/Bilderkennung/Produktspezifikation/JetsonNano/Inhalt/JetsonDigits6.tex $
% $Version: 1765 $
%
%%%%%%%%%%%%



\chapter{DIGITS}


{\tiny Source: \href{https://github.com/NVIDIA/nvidia-docker/wiki/DIGITS}{Jonathan Calmels edited this page Nov 16, 2017 ·  14 revisions}}


\section{Description}

Docker image based on \href{https://github.com/NVIDIA/DIGITS}{DIGITS}.

\section{Requirements}

See the requirements for \href{https://github.com/NVIDIA/nvidia-docker/wiki/CUDA#requirements}{CUDA} since the DIGITS image is based on a cuDNN runtime image.

\section{Examples}

\begin{lstlisting}
# Run DIGITS on host port 5000
docker run --runtime=nvidia --name digits -d -p 5000:5000 nvidia/digits
\end{lstlisting}

If you want to use a dataset stored in a host directory, you will need to import it inside the container using a volume:

\begin{lstlisting}
# Run DIGITS with the mnist dataset stored on the host under /opt/mnist
docker run --runtime=nvidia --name digits -d -p 5000:5000 -v /opt/mnist:/data/mnist nvidia/digits
\end{lstlisting}

If you want to persist jobs across multiple DIGITS containers, you can use a named volume. For DIGITS 3.0:

\begin{lstlisting}
# Run DIGITS storing jobs in a host volume named digits-jobs

# For DIGITS 3.0
docker run --runtime=nvidia --name digits -d -p 5000:34448 -v digits-jobs:/usr/share/digits/digits/jobs nvidia/digits:3.0

# For DIGITS 3.3 and 4.0
docker run --runtime=nvidia --name digits -d -p 5000:34448 -v digits-jobs:/jobs nvidia/digits:4.0

# For DIGITS 5.0
docker run --runtime=nvidia --name digits -d -p 5000:5000 -v digits-jobs:/jobs nvidia/digits:5.0

# For DIGITS 6.0
docker run --runtime=nvidia --name digits -d -p 5000:5000 -p 6006:6006 -v digits-jobs:/jobs nvidia/digits:6.0
\end{lstlisting}

\section{Tags available}

Check the \href{https://hub.docker.com/r/nvidia/digits/}{DockerHub}
